%% kapa: exam=1 round=28 year=2017 period=1 form=A subject=economics subject_ko=경제학원론 seq=18 number=58
%% source: Q-Net 감정평가사 기출문제 (cst00309) artlSeq=5199603 file=extracted/1008/5199603/2017년 제28회 감정평가사 1차시험 문제지/1교시 A형.pdf
%% conversion: kordoc 4.24.51 -> tools/gen1.py (auto); review=proofread-sonnet; flags=none
\begin{question}{58}
정보의 비대칭성에 관한 설명으로 옳지 않은 것은?
\begin{choices}
  \item 사고가 발생할 가능성이 높은 사람일수록 보험에 가입할 가능성이 크다는 것은 역선택(adverse selection)에 해당한다.
  \item 화재보험 가입자가 화재예방 노력을 게을리 할 가능성이 크다는 것은 도덕적 해이(moral hazard)에 해당한다.
  \item 통합균형(pooling equilibrium)에서는 서로 다른 선호체계를 갖고 있는 경제주체들이 동일한 전략을 선택한다.
  \item 선별(screening)은 정보를 보유하지 못한 측이 역선택 문제를 해결하기 위해 사용할 수 있는 방법이다.
  \item 항공사가 서로 다른 유형의 소비자에게 각각 다른 요금을 부과하는 행위는 신호 발송(signaling)에 해당한다.
\end{choices}
\end{question}
\answer{5}
